\section{Introduction}
\label{sec:introduction}

Language agents are increasingly used as interactive software-engineering assistants, with repository-level systems and benchmarks exposing agents to real issues, codebases, and executable tests \cite{jimenez2024swebench,yang2024sweagent}. Their ability to improve from prior debugging experience is less settled. A debugging agent may observe many traces of failure reproduction, code inspection, patch attempts, and test reruns. The central question is whether those traces can be compressed into a form that helps the agent solve new, related tasks without leaking task-specific fixes or transferring the wrong lesson.

Existing memory mechanisms represent experience as verbal reflections, distilled insights, conditional guidelines, instruction manuals, retrieved workflows, or procedural repositories \cite{shinn2023reflexion,zhao2024expel,fu2024autoguide,chen2024automanual,wang2025workflowmemory,fang2026memp}. These systems establish that agent experience can be stored and reused, but they differ in whether they specify applicability, ordered action, and conditions for rejecting prior experience. This distinction is important in debugging. Two tasks can share a surface failure family while differing in the invariant that must be repaired. A memory that is useful on one task can therefore become harmful on a nearby held-out task if it focuses the agent on the wrong formatter behavior, file region, or patch pattern.

Recent memory systems already report process measures such as actions, steps, and token consumption, and they recognize that stale or noisy memories can reduce performance \cite{fang2026memp,ouyang2026reasoningbank}. The unresolved evaluation question is narrower: whether a fixed procedure induced from a few debugging trajectories transfers to held-out bugs, how it changes the repair process, and whether observed failures can be attributed to memory misapplication. Final solve rate alone cannot distinguish procedural transfer from added context or from an adjacent but wrong repair pattern. We therefore measure diagnosis-edit-test cycles, token cost, failed patches, repeated mistakes, and explicit negative transfer alongside repair success.

This paper studies a more procedural form of agent memory. We define a \texttt{SKILL.md} as a \emph{trajectory-induced procedural memory artifact}: a natural-language artifact induced from a small number of prior debugging trajectories, scoped to a task family, and constrained to encode three reusable components. First, it states trigger conditions describing when the procedure should be considered. Second, it provides an ordered debugging procedure for reproducing the failure, localizing the invariant, editing minimally, and validating the repair. Third, it lists failure modes that should stop, redirect, or qualify the procedure. The artifact is natural language because the downstream agent can execute natural-language procedures, not because the contribution is a manually tuned prompt.

We call this artifact procedural memory because it is produced by an induction pipeline from prior trajectories, constrained by leakage rules, deployed unchanged on held-out tasks, and evaluated by transfer behavior rather than by manual prompt quality.

Our main contribution is an evaluation protocol for low-shot procedural skill transfer in language-agent debugging. The protocol asks whether an agent can induce a \texttt{SKILL.md} from a small set of training trajectories and then use it on within-family held-out debugging tasks. It evaluates memory artifacts by solve rate, recorded diagnosis-edit-test cycles, tokens per solved task, failed patches, repeated mistakes, and negative transfer. This protocol is designed to separate three questions that are often conflated: whether memory improves final repair success, whether it improves the repair process, and whether it can misdirect the agent when a learned procedure is applied outside its true scope.

The empirical setting is intentionally controlled. We construct a hard-smoke slice of verified PyBugHive \texttt{black} formatter tasks after easier pilot tasks failed to separate memory methods \cite{antal2024pybughive}. Formatter tasks are useful for this first study because they expose localized but nontrivial behavioral invariants, along with adjacent subfamily failures where the wrong procedure can plausibly transfer. The primary study uses six held-out tasks run under the same model and harness configuration. All methods receive the same task prompt, isolated workspace, test command, budget, and leakage constraints. We compare Auto SKILL.md against no memory, a generic debugging checklist, Reflexion-style memory \cite{shinn2023reflexion}, and a length-matched Reflexion control. The length-matched control is especially important because it tests whether improvements come merely from adding more memory text or from organizing experience as trigger conditions, procedure, and failure modes.

In the primary single-execution run on this six-task hard-smoke slice, Auto SKILL.md matches the best observed solve rate at 5/6. Reflexion memory and the generic checklist also solve 5/6 tasks, and the checklist matches Auto SKILL.md at 2.0 cycles per solved task. Auto SKILL.md uses fewer mean tokens, but task-level uncertainty around the checklist comparison is wide. Its clearest observed process differences are against reflection baselines on common-solved tasks. At the same time, it produces the only explicit negative-transfer flag in the primary run. This failure is not incidental to the paper's contribution; it shows why procedural memory should be evaluated for misapplication rather than reported only as an aggregate success rate.

We also include a supplementary structural-control rerun under \texttt{gpt-5.5}, comparing Auto SKILL.md with a format-shuffled SKILL.md that preserves the content of the induced skill while disrupting its section organization. Both methods solve all six tasks. Auto SKILL.md uses fewer recorded cycles and fewer failed patches, while the shuffled control uses fewer tokens in aggregate. This result narrows the structural claim. The evidence does not show that the canonical three-section SKILL.md format is necessary for solvability under a stronger model. Instead, it supports a more cautious process-behavior claim: procedural organization changes how the agent searches, edits, and validates, and the proposed protocol makes those differences measurable.

This framing also clarifies the paper's scope. We do not claim that SKILL.md universally outperforms Reflexion memory, that six tasks establish statistical significance, or that formatter-style debugging generalizes to all software-engineering agents. The present study is a controlled first step. It contributes a task and trajectory schema, an induction protocol, memory baselines, process metrics, and negative-transfer accounting for evaluating low-shot procedural skill transfer. The broader goal is to make skill self-evolution empirically testable: a self-improving agent should not only accumulate memories, but also learn when those memories apply, how they should guide action, and when they should be distrusted.

The paper makes three contributions:
\begin{enumerate}
\item We formalize natural-language \texttt{SKILL.md} as a trajectory-induced procedural memory artifact for language-agent debugging, distinct from unstructured reflection memory and generic debugging advice.
\item We introduce a reproducible evaluation protocol for low-shot within-family procedural skill transfer, including task schemas, trajectory schemas, memory controls, process-efficiency metrics, and negative-transfer annotation.
\item We report a controlled PyBugHive \texttt{black} hard-smoke pilot in which Auto SKILL.md matches the best observed solve rate and shows initial task-level process differences from unstructured memory controls, together with a supplementary format-shuffled control that bounds the structural claim.
\end{enumerate}
